In the last section has been discussed the method with which the \gls{IP}s are found.
In this section the actual instrumentation is reviewed.
As stated in \cref{sec:instr_point_search} there are two types of instrumentation :
\begin{itemize}
    \item The first one is used to setup the variables used by the \gls{SHM} update rule.
    \item All the other \gls{IP} are used for the actual \gls{SHM} update.
\end{itemize}
With reference to \cref{alg:afl_style_update} there are 3 main variables to consider:
\begin{itemize}
    \item \textit{current\_bb\_id} : ID of the \gls{BB} in which the instrumentation is placed
    \item \textit{previous\_bb\_id} : ID of the \gls{BB} from which the branch is taken
    \item \textit{shared\_map} : Address of the \gls{SHM}
\end{itemize}
In the last example of the previous section (\cref{fig:pow_example}) the \gls{IP}s are already been found.
Without considering the first (that is where the initial instrumeentation is placed), in all the other \gls{IP}s the 
initial address of the \gls{BB} is where the instrumentation will be placed.
Lets consider for example \textbf{0x11f1} \gls{IP} with previous \gls{BB} equal \textit{0x11c2} and the current \textit{0x11f1}.
In this case, since there is only one incoming edge to \textbf{0x11f1}, the \textit{previous\_bb\_id} could be placed inline.
But in the general case (e.g. \textbf{0x11b5}) there are more edges that ends in the \gls{BB}. For this reason become impossible 
place the ID of the \gls{BB} directly inline.
The solution implemented in the tool discussed in this thesis consist in storing a variable in a section of the binary\footnote{
    This section could be the same used for relocation or a new one added specifically to store the data
}
This variable represent the \textit{previous\_bb\_id} and it is initialized by the initial instrumentation to a random value.

For the \textit{current\_bb\_id} instead there is no need to store it into memory. It will be hardcoded into the 
assembly instrumented within the basic block (to a random value).

The \textit{shared\_map} is already been resolved and written in the section added to the binary specifically to store 
resolved symbols.

So now all the instrumentation phase could be revised with greter details.
\begin{center}
    \textbf{Initial Instrumentation}
\end{center}
This instrumentation is needed to :
\begin{itemize}
    \item Initialize the \textit{previous\_bb\_id}. This value is stored into a specific section created for the porpouse of 
        holding temporary values needed for instrumentation.
    \item Dereference the resolved \gls{SHM} address. This initialization step is performed to optimization porpouse. 
        As can be seen in \cref{fig:shm_relocation} the instrumented section that work as \gls{GOT} holds inside the address 
        of the variable inside the injected \gls{SO}. In order to obtain the actual \gls{SHM} address the value must be dereferenced twice. 
        If, at any \gls{IP}, a dereference step must be performed twice, this will result in many useless overheads.
        For this reason in the initial instrumentation the value is dereference twice and then saved into the same section of \textit{previous\_bb\_id}.
        By doing this, at each \gls{IP}, the \gls{SHM} address must only be loaded into a register and find the offset of the correct cell.
\end{itemize}
\begin{center}
    \textbf{\gls{SHM} update instrumentation}
\end{center}
In \cref{alg:shm_update} can be seen the general approach to \gls{SHM} update.
Since the \texttt{rax} register value is modified, its value must be saved and restore before the actual \gls{SHM} update.
The value of the \texttt{rax} is not saved into the stack but in the same section where \textit{previous\_bb\_id} and 
\textit{shared\_map\_address} are stored.
\begin{algorithm}
    \caption{Pseudo-code of Shared Memory update instrumentation}
    \label{alg:shm_update}
    \begin{algorithmic}[1]
        \STATE $\text{Save the state of \texttt{rax} register that will be used for computation}$
        \STATE $\texttt{rax} \gets \textit{previous\_bb\_id}$
        \STATE $\texttt{rax} \gets \texttt{rax} \oplus \textit{current\_bb\_id}$ \COMMENT{\small \textcolor{green!60!black}{into \texttt{rax} there will be the result of hash operation between the basic blocks}}
        \STATE $\texttt{rax} \gets \texttt{rax} + \textit{shared\_map\_address} $\COMMENT{\small \textcolor{green!60!black}{since the shared map is an array, the offset is obtained adding the initial address plus the index}}
        \STATE $[\,\texttt{rax}\,] \gets [\,\texttt{rax}\,] + 1$ \COMMENT{\small \textcolor{green!60!black}{the [\,\,] is the dereference operator}}
        \STATE $\textit{previous\_bb\_id} \gets \textit{current\_bb\_id} \gg 1$
        \STATE $\text{Restore the value of \texttt{rax} register to the previously saved value}$
    \end{algorithmic}
\end{algorithm}

Now the complete approach could be revised : 
\begin{itemize}
    \item Add a dynamic dependency to a \gls{SO} that contains the forkserver function.
    \item Add a symbol that will contains the \gls{SHM} address.
    \item Add two new sections. The first will be used as \gls{GOT} to relocate the symbol. The second will 
        be used as temporary memory space to store \textit{previous\_bb\_id} and \textit{shared\_map\_address} alongside register saving.
    \item Add a relocation entry in order to resolve the symbol previously added. That symbol is exported from the \gls{SO}.
    \item Find the \gls{IP}s with the algorithm discussed in \cref{sec:instr_point_search}
    \item Add the initial instrumentation used to initialize the values of 
        \textit{previous\_bb\_id} and \textit{shared\_map\_address} stored in the added section
    \item For each previously found \gls{IP} add the \gls{SHM} update instructions as seen in \cref{alg:shm_update}
\end{itemize}

\begin{center}
    \textbf{Optimization}
\end{center}
Since fuzzing is a technique that intensively rely on execution speed with the single input, the optimization 
of the instrumentation could save time.
In the previously discussed approach there is a major cause of overheads, that is the point in which the forkserver is executed.
Infact, as can be seen in \cref{fig:dynamicallyLinkedProgramLoad}, since the forkserver is executed by the dynamic linker 
before the execution of the actual target \texttt{\_start} function, at each input the target binary must execute many operations 
before actually reach the main:
\begin{itemize}
    \item \_start function must be executed
    \item \_\_libc\_start\_main function must be executed, in which the environment is setted up for the main execution
    \item .init\_array functions are scanned and eventually executed
\end{itemize}
The ideal approach is to execute the forkserver at the very beginning of the main function, 
avoiding execute all the preamble functions for each input.
To achieve this the relocation mechanism is used again.
In the injected \gls{SO}, the constructor function does not contains the forkserver anymore, but is used only to 
prepare the execution of the forkserver. The workflow can be seen in \cref{fig:final_forkserver_invocation}
Precisely the constructor function will :
\begin{itemize}
    \item setup the \gls{SHM} reading from the environmental variable the \gls{SHM} ID and assign to the exported symbol its address.\footnote{
        In the next chapter the \gls{SHM} attach will be discussed deeply. Infact the \gls{SO} also includes a \textit{standalone-mode} 
        in which there is no need for a \gls{SHM}
    }
    \item assign to a symbol (that will be also exported) the address of the forkserver function that will be executed
\end{itemize}
To the target binary will be added two symbols instead of one :
\begin{itemize}
    \item The first is the usual \gls{SHM} address
    \item The second is a pointer to the forkserver function
\end{itemize}
These two symbols are then resolved by the dynamic linker at binary startup and their value is written to the section that 
is has been added.
Into the initial instrumentation first of all the address of the forkserver function is dereferenced and then executed.
At the exit of forkserver the rest of the main function is executed without reexecuting all the binary startup functions 
for each input.
\begin{figure}[ht]
    \centering
    \includegraphics[width=400pt]{FinalForkserverExecutionDiagram.png}
    \caption{Optimized forkserver invocation}
    \label{fig:final_forkserver_invocation}
\end{figure}